% Generated draft: requires mathematical and source review.
\documentclass[11pt]{article}
\usepackage[margin=25mm]{geometry}
\usepackage{fontspec}
\setmainfont{Latin Modern Roman}
\usepackage{amsmath,amssymb,mathtools,unicode-math}
\setmathfont{Latin Modern Math}
\usepackage{enumitem,needspace,xurl,hyperref}
\setlength{\parindent}{0pt}
\setlength{\parskip}{5pt}
\setlength{\emergencystretch}{3em}
\newcommand{\N}{\mathbb{N}}
\newcommand{\Z}{\mathbb{Z}}
\newcommand{\Q}{\mathbb{Q}}
\newcommand{\R}{\mathbb{R}}
\newcommand{\C}{\mathbb{C}}
\newcommand{\F}{\mathbb{F}}
\newcommand{\A}{\mathbb{A}}
\newcommand{\PP}{\mathbb{P}}
\newcommand{\E}{\mathbb{E}}
\newcommand{\eps}{\varepsilon}
\newcommand{\dd}{\,\mathrm{d}}
\begin{document}
\section*{UnsolvedMath ID: 700939}
\textbf{Source problem code:} catalogue entry
\subsection{Definitions and mathematical statement}
Let \(p\) be a prime. A \(p\)-group here means a group in which every element has finite order equal to a power of \(p\). A group \(G\) is residually finite if, for every \(g\ne 1\), there is a homomorphism from \(G\) to a finite group whose image of \(g\) is nonidentity. Let \(\mathcal C_p\) denote the class of infinite finitely generated residually finite \(p\)-groups.

For \(G\in\mathcal C_p\), fix a finite generating set \(S\). With the usual convention of allowing generators and their inverses, define
\[
\delta_S(g)=\min\{k\geq 0:g=s_1\cdots s_k,\ s_i\in S\cup S^{-1}\},
\qquad
\operatorname{ord}_G(g)=\min\{m\geq 1:g^m=1\}.
\]
For integers \(n\geq 0\), put
\[
\pi_{G,S}(n)=\max\{\operatorname{ord}_G(g):g\in G,\ \delta_S(g)\leq n\}.
\]
This maximum exists because each word ball is finite. The notation \(|g|\) in the source denotes \(\operatorname{ord}_G(g)\).

\textbf{Part (b), the remaining question in the supplied record.} What is the minimal possible asymptotic rate of growth of \(\pi_{G,S}(n)\) for the class of groups indicated in part (a)? The underlying group hypotheses in that part are \(G\in\mathcal C_p\). However, the reference to that class does not explicitly resolve whether the additional subpower bound asked for in part (a) is itself a restriction in part (b). The source also leaves the comparison convention for growth rates and the quantification over \(p\) unspecified; no precise extremal ordering or assertion of an attained minimum is imposed here.

\textbf{Part (a), retained as archived context.} Does there exist \(G\in\mathcal C_p\), a finite generating set \(S\), and a function \(f:\{0,1,2,\ldots\}\to[1,\infty)\) such that
\[
\forall g\in G,\qquad \operatorname{ord}_G(g)\leq f(\delta_S(g)),
\qquad
\forall\lambda>0,\qquad
\lim_{n\to\infty}\frac{f(n)}{n^\lambda}=0?
\]
The domain and value at zero merely make the bound include the identity. This is the subpower interpretation of the source's phrase that \(f\) grows more slowly than every positive power; no monotonicity condition is supplied. The source does not specify whether \(p\) is prescribed or may be chosen.

The supplied archive reports an affirmative answer to part (a), attributed to A. V. Rozhkov. As background to part (b), the source attributes to R. I. Grigorchuk the existence of \(p\)-groups with a bound \(\pi(n)\leq n^\lambda\) for some \(\lambda>0\), and states that a result of E. I. Zel'manov implies that \(\pi(n)\) is unbounded when \(G\) is infinite in the indicated setting. These are supplied historical assertions, not independently verified results in this draft.
\subsection{Short English statement}
For an infinite finitely generated residually finite \(p\)-group, measure how large an element order can occur among words of length at most \(n\). Part (b) asks how slowly this maximum can grow. The supplied archive reports that part (a), asking for growth below every positive power, was answered affirmatively. The precise growth comparison, prime quantification, and scope of the class in part (b) require clarification.
\subsection{Sources}
\begin{itemize}
\item Local UnsolvedMath registry, supplied complete record with permanent ID \(700939\), exact source code \texttt{catalogue entry}, title \emph{Kourovka Notebook Problem 13.21}. The supplied status is \texttt{partially\_solved}; it records archived solved part (a) and remaining part (b).
\item \emph{The Kourovka Notebook (Unsolved Problems in Group Theory)}, 13th issue (1995), Problem 13.21, proposed by R. I. Grigorchuk; supplied background identifies an update in the 21st issue (2026), \texttt{arXiv:1401.0300v48}. The supplied collection description names E. I. Khukhro and V. D. Mazurov as editors.
\item Archive of Solved Problems: affirmative answer to part (a), attributed to A. V. Rozhkov, \emph{Dokl. Math.}, \textbf{58}, no. 2 (1998), 234--237. No article title or link is supplied.
\item R. I. Grigorchuk, \emph{Math. USSR--Izv.}, \textbf{25} (1985), 259--300, cited in the source for a polynomial upper bound on \(\pi\) in examples of \(p\)-groups. No article title or link is supplied.
\item E. I. Zel'manov: the supplied statement attributes unboundedness of \(\pi\) for infinite groups in the indicated setting to a result of his, without identifying a publication.
\item The supplied status review is dated \texttt{2026-10-16}, with evidence recorded as \texttt{published}. Its assessment that archived part(s) are solved and remaining part(s) appear in the main section is supplied dated curation, not independently verified here.
\end{itemize}
\end{document}
